% Generated draft: requires mathematical and source review.
\documentclass[11pt]{article}
\usepackage[margin=25mm]{geometry}
\usepackage{fontspec}
\setmainfont{Latin Modern Roman}
\usepackage{amsmath,amssymb,mathtools,unicode-math}
\setmathfont{Latin Modern Math}
\usepackage{enumitem,needspace,xurl,hyperref}
\setlength{\parindent}{0pt}
\setlength{\parskip}{5pt}
\setlength{\emergencystretch}{3em}
\newcommand{\N}{\mathbb{N}}
\newcommand{\Z}{\mathbb{Z}}
\newcommand{\Q}{\mathbb{Q}}
\newcommand{\R}{\mathbb{R}}
\newcommand{\C}{\mathbb{C}}
\newcommand{\F}{\mathbb{F}}
\newcommand{\A}{\mathbb{A}}
\newcommand{\PP}{\mathbb{P}}
\newcommand{\E}{\mathbb{E}}
\newcommand{\eps}{\varepsilon}
\newcommand{\dd}{\,\mathrm{d}}
\begin{document}
\section*{UnsolvedMath ID: 6800020}
\textbf{Source problem code:} catalogue entry
\subsection{Definitions and mathematical statement}
The supplied entry contains three distinct problems.

\textbf{Normal holonomy.} Let $M^n\subset S^N$, $n\geq2$, be an irreducible full homogeneous submanifold. Full means not contained in a proper totally geodesic subsphere. Its normal connection in the sphere has a holonomy action on a normal fiber. If this action is not transitive on the unit normal sphere, must $M$ be, up to congruence, an orbit of an $s$-representation? An $s$-representation is the isotropy representation of a Riemannian symmetric space. The precise irreducibility convention and whether full or restricted normal holonomy is intended are not supplied. Normal holonomy in the sphere is distinguished from the Euclidean normal bundle, which includes a radial direction.

\textbf{Index.} For an irreducible Riemannian symmetric space $S$, let $i(S)$ be the minimum codimension of a proper complete connected totally geodesic submanifold. Let $i_r(S)$ be the minimum codimension among reflective submanifolds, meaning connected components of fixed-point sets of involutive isometries. Prove
\[
i(S)=i_r(S)
\]
for the source's class excluding $G_2/SO(4)$ and its symmetric dual. The excluded spaces are not added back as part of this assertion despite the summary's reference to them.

\textbf{Nullity.} For a Riemannian manifold define
\[
\mathcal N_p=\{v\in T_pM:R_p(v,w)z=0\text{ for all }w,z\in T_pM\}.
\]
Homogeneity makes its dimension constant. For homogeneous manifolds with nontrivial nullity, the source asks separately: are there examples that are not topologically trivial; examples $M=G/H$ with $G$ nonsolvable; K\"ahler examples; and examples in every dimension $d\geq5$? A K\"ahler manifold has a metric-compatible parallel complex structure and therefore even real dimension; the K\"ahler question is separate from the every-dimension question.

The intended nullity class is incomplete. Flat manifolds and products with Euclidean factors already have nontrivial nullity, so nonflatness, irreducibility, absence of Euclidean factors, or nonparallel nullity may be intended but are not stated. The meaning of topologically trivial is also missing. The $G/H$ question needs an effectiveness convention, since adding a group factor acting trivially would alter the solvability of a presenting group without changing the manifold.
\subsection{Short English statement}
Resolve the normal-holonomy orbit conjecture, the index/reflective-index equality outside the named exceptions, and four separate existence questions for homogeneous manifolds with nontrivial nullity. Clarify the nullity class and group-action conventions.
\subsection{Sources}
\begin{itemize}
\item Morgan and Pansu, \emph{A List of Open Problems in Differential Geometry} (2018), Question environment 20, ``Homogeneous Riemannian manifolds with nontrivial nullity,'' proposed by Carlos Olmos, C\'ordoba, Argentina, \url{https://www.imo.universite-paris-saclay.fr/~pansu/problems_MTDG.tex}; accessed 2026-07-29, source-TeX extraction.
\item The supplied research summary mentions Olmos and collaborators' partial results without individual publication citations.
\item Local registry \texttt{data/unsolvedmath/problems.json}, permanent UnsolvedMath ID $6800020$, source code catalogue entry, same titled problem. Dataset attribution: Ulam AI and the UnsolvedMath Contributors.
\end{itemize}
\end{document}
